\begin{frame}{2차원 GDA 시각화}
  \begin{columns}[T]
    \begin{column}{0.55\textwidth}
      \begin{itemize}
        \item 2차원 가우시안 300+300개 샘플로 검증:
          참 평균 $[1,2]$, $[4,5]$, 참 공분산
          $\begin{bmatrix}1.5&0.8\\0.8&1.2\end{bmatrix}$.
        \item 추정 $\hat\mu_0,\hat\mu_1$ 이 참 평균에
          \textbf{접근} (300개 표본이라 노이즈 $\pm0.1$ 정도).
        \item 결정 경계 = $w_0 x_1 + w_1 x_2 + b = 0$ 인 \textbf{직선}.
        \item 기울기는 공분산의 ``퍼진 방향''에 맞춰 보정된 $w$ 가
          결정.
      \end{itemize}
    \end{column}
    \begin{column}{0.43\textwidth}
      \centering
      \includegraphics[width=\linewidth]{ch03_2_gda_boundary.png}
      \\[0.3em]{\footnotesize MLE로 추정한 두 클래스 가우시안과
      결정 경계 직선(검은 점선).}
    \end{column}
  \end{columns}
\end{frame}
